\twocolumn[{
\section{Discussion and Reproducibility}
\begin{center}\small
\begin{tabular}{p{.25\textwidth}p{.42\textwidth}p{.24\textwidth}}
\toprule
Additional experiment family & Main evidence & Decision\\
\midrule
Graph~\cite{graph} / soft or selective pseudo labels & Graph validation F1 0.8932; public F1 0.84081 versus incumbent 0.85405. & Negative generalization example.\\
Primary/secondary reasons and whitening & Primary/secondary OOF 0.8960, validation 0.8590; shrinkage whitening did not pass checks. & Retain summary.\\
Label shift~\cite{saerens,bbse}; residual Other/text & Prior correction gains one late trial; residual Other worsens late-year F1. & No promotion.\\
Local ensemble, Other kernel, F1 decoder & Fixed-ridge OOF 0.9026; sequential procedure 0.9002. & Retain public 0.88568.\\
Ridge grid and group weighting & Grid-selected OOF 0.8983; weighting changes no final OOF decisions. & Retain public 0.88568.\\
\bottomrule
\end{tabular}
\captionof{table}{Later unsuccessful hypotheses. Earlier 157-label OOF protocols differ from
the newer 235-label near-group protocol and cannot be ranked directly. Fixed-recipe studies
followed recipe selection. Sources and recorded evidence are included in the code ZIP.}
\label{tab:negative}
\end{center}
}]
\subsection{Why more features need not help Part 1}
The embedding-only choice matches the target: stopping-reason text can express recruitment,
funding, or administrative problems that a drug's biology does not describe.
The pathway decline in Section~\ref{sec:p1features} supports excluding that tested block.
Possible explanations are incomplete exact-name drug matching (about 56\% coverage in Part 1~\cite{handout}),
coarse or redundant pathways, and extra coefficients estimated from only 157 gold labels;
pseudo targets add uncertain supervision, not 1,014 independent gold labels.
Missing drug matches also do not imply biological absence. These are plausible mechanisms,
not separately established causes of the observed decline.

Criteria, disease, drug, ICD, target/action, and molecular representations can increase variance
or dilute reason similarity. No exhaustive all-feature or fingerprint ablation was run;
their usefulness is unresolved.
For the chosen kernel, global scaling preserves Euclidean geometry while column-wise scaling
alters it. The public improvement also changes model family and gold count (157 to 235),
so it cannot isolate geometry. Local models avoid inheriting uncertain cluster targets;
rare-label suppression is a limitation of the tested K-Means/GMM fits, not all semi-supervised methods.

Follow-ups did not improve the pipeline. The sample-F1 decoder assumes independent Bernoulli
marginals, potentially violated by correlated reasons. A 16-setting ridge grid and inverse-group
weighting also failed: 82/235 labels occupy 13 repeated groups, mostly Recruitment,
but downweighting did not improve held-group decisions.

\subsection{Selection and temporal limitations}
Validation labels and public feedback were inspected repeatedly. Later nested resampling
controls leakage within each fit but cannot erase that search history. The final Part 1 fit
uses all 157+78 labels after recipe selection. Historical validation scores come from separate
157-label fits; they are not independent scores of the final 235-label artifact.

The late-20 and late-56 diagnostics have different training populations and reason mixtures:
late train contains one Recruitment-positive trial, versus 21 in late validation. Neither
small subset determines future performance. Temporal fits retain earlier duplicates available at that date; saved diagnostics
isolate novel groups. Test data never enter fitting or selection. Private performance is unknown.

\subsection{Reproduction and preserved evidence}
The report uses CVPR style~\cite{template}. \path{models/selection.json} fixes final hashes.
Run \texttt{python reproduce.py} from the project root to verify saved predictions and IDs.
The ZIP includes data/audits, models/evidence, source/tests/notebooks, and report/LaTeX in
\texttt{data/}, \texttt{models/}, \texttt{tools/}, and \texttt{report/}, respectively.
Run \path{report/scripts/report_analysis.py} then \path{report/scripts/generate_tables.py}
with Python to regenerate metrics, figures, and tables. Five notebooks in \path{tools/notebooks/}
replay saved evidence with training disabled. Historical runners/manifests are included;
full backups and the frozen November 1 delivery remain local.

\paragraph{Team contributions and assistance.}
Pei-Chen Chung (pc786) and Yu-Jye Hu (yh2444) contributed equally to data preprocessing,
exploratory data analysis, feature engineering, model development, performance evaluation,
and preparation of the final report. Both members shared responsibilities, jointly discussed
key modeling decisions, and contributed equally to the final deliverables. OpenAI Codex
assisted with implementation, verification, organization, and drafting. Both authors are
responsible for the submitted work.
